\documentclass[11pt,a4paper]{article}

\usepackage[utf8]{inputenc}
\usepackage[T1]{fontenc}
\usepackage{lmodern}
\usepackage[margin=1in]{geometry}
\usepackage{hyperref}
\usepackage{xcolor}
\usepackage{listings}
\usepackage{enumitem}
\usepackage{graphicx}
\usepackage{titlesec}
\usepackage{tcolorbox}
\usepackage{array}
\usepackage{booktabs}
\usepackage{longtable}

% ---------- Hyperref ----------
\hypersetup{
    colorlinks=true,
    linkcolor=blue!60!black,
    urlcolor=blue!60!black,
    citecolor=blue!60!black,
    pdftitle={Deploying ML Models on AWS SageMaker},
    pdfauthor={Prabal}
}

% ---------- Listings (code) ----------
\definecolor{codebg}{RGB}{245,245,245}
\definecolor{codekw}{RGB}{0,0,180}
\definecolor{codestr}{RGB}{163,21,21}
\definecolor{codecom}{RGB}{0,128,0}

\lstdefinestyle{pycode}{
    backgroundcolor=\color{codebg},
    basicstyle=\ttfamily\footnotesize,
    keywordstyle=\color{codekw}\bfseries,
    stringstyle=\color{codestr},
    commentstyle=\color{codecom}\itshape,
    language=Python,
    breaklines=true,
    showstringspaces=false,
    numbers=left,
    numberstyle=\tiny\color{gray},
    frame=single,
    rulecolor=\color{gray!30},
    tabsize=2,
    captionpos=b
}

\lstdefinestyle{bashcode}{
    backgroundcolor=\color{codebg},
    basicstyle=\ttfamily\footnotesize,
    keywordstyle=\color{codekw}\bfseries,
    stringstyle=\color{codestr},
    commentstyle=\color{codecom}\itshape,
    language=bash,
    breaklines=true,
    showstringspaces=false,
    frame=single,
    rulecolor=\color{gray!30},
    tabsize=2
}

% ---------- Note boxes ----------
\newtcolorbox{notebox}{
    colback=blue!5!white,
    colframe=blue!50!black,
    title=Note,
    fonttitle=\bfseries,
    sharp corners,
    boxrule=0.5pt
}

\newtcolorbox{warnbox}{
    colback=orange!5!white,
    colframe=orange!70!black,
    title=Important,
    fonttitle=\bfseries,
    sharp corners,
    boxrule=0.5pt
}

\newtcolorbox{tipbox}{
    colback=green!5!white,
    colframe=green!50!black,
    title=Tip,
    fonttitle=\bfseries,
    sharp corners,
    boxrule=0.5pt
}

% ---------- Titles ----------
\titleformat{\section}{\Large\bfseries\color{blue!50!black}}{\thesection}{1em}{}
\titleformat{\subsection}{\large\bfseries\color{blue!40!black}}{\thesubsection}{1em}{}

\title{\textbf{Deploying Machine Learning Models on AWS SageMaker}\\
\large A Step-by-Step Deployment Guide}
\author{Prabal}
\date{\today}

\begin{document}

\maketitle
\tableofcontents
\newpage

% =========================================================
\section{Introduction}
% =========================================================

Amazon SageMaker is AWS's fully managed machine learning service that covers the
entire ML lifecycle: data labeling, training, tuning, deployment, and monitoring.
For deployment specifically, SageMaker abstracts away the heavy lifting of
provisioning compute, setting up HTTPS endpoints, load balancing, autoscaling,
and rolling updates, letting teams ship production models with a few API calls.

This document focuses on the \textbf{deployment} phase, assuming a model has
already been trained (in SageMaker, locally, or elsewhere). It walks through the
end-to-end flow of exposing that model as a production inference service.

\begin{notebox}
SageMaker offers four deployment modes:
\begin{itemize}[leftmargin=*,nosep]
    \item \textbf{Real-time inference} --- persistent HTTPS endpoint, low latency.
    \item \textbf{Serverless inference} --- pay-per-request, scales to zero.
    \item \textbf{Asynchronous inference} --- queued requests, large payloads, long timeouts.
    \item \textbf{Batch transform} --- offline bulk scoring over S3 datasets.
\end{itemize}
The main flow below targets \emph{real-time inference}, which is the most common mode.
\end{notebox}

% =========================================================
\section{Core Concepts}
% =========================================================

Before deploying, understand the three SageMaker primitives that compose an endpoint:

\begin{longtable}{p{3.2cm} p{11.5cm}}
\toprule
\textbf{Primitive} & \textbf{Purpose} \\
\midrule
\texttt{Model} & A logical object that binds a \emph{model artifact} (weights stored
in S3 as \texttt{model.tar.gz}) to an \emph{inference container image} (from
ECR) and an \emph{IAM role}. \\
\addlinespace
\texttt{EndpointConfig} & Declares one or more \emph{production variants} --- which
model(s) to serve, instance type, instance count, autoscaling policy, and traffic
split. \\
\addlinespace
\texttt{Endpoint} & The live, invokable HTTPS resource. Built from an
\texttt{EndpointConfig}. Updating it performs a safe blue/green rollout. \\
\bottomrule
\end{longtable}

\begin{tipbox}
Keep these three objects separate in your mind: \textbf{what to serve} (Model),
\textbf{how to serve it} (EndpointConfig), and \textbf{where it is served}
(Endpoint). Blue/green deployments and A/B traffic splits all happen by creating
new EndpointConfigs and updating the Endpoint to point at them.
\end{tipbox}

% =========================================================
\section{Prerequisites}
% =========================================================

\begin{enumerate}[leftmargin=*]
    \item An AWS account with billing enabled.
    \item AWS CLI installed and configured (\texttt{aws configure}).
    \item Python 3.9+ with \texttt{boto3} and \texttt{sagemaker} SDK installed:
    \begin{lstlisting}[style=bashcode]
pip install boto3 sagemaker
    \end{lstlisting}
    \item An IAM execution role for SageMaker with the \texttt{AmazonSageMakerFullAccess}
          managed policy, plus S3 read access to your model bucket. The role's
          trust policy must allow \texttt{sagemaker.amazonaws.com} to assume it.
    \item An S3 bucket in the same region as your intended endpoint.
    \item A trained model (e.g., a PyTorch \texttt{.pt} file, TensorFlow
          \texttt{SavedModel}, scikit-learn pickle, or HuggingFace checkpoint).
\end{enumerate}

\begin{warnbox}
Region consistency matters. The IAM role, S3 bucket, ECR image, and SageMaker
endpoint must all be in the \emph{same AWS region}. Cross-region references
cause opaque permission errors and latency overhead.
\end{warnbox}

% =========================================================
\section{Step-by-Step Deployment}
% =========================================================

% ---------------------------------
\subsection{Step 1 --- Package the Model Artifact}
% ---------------------------------

SageMaker expects the model artifact as a gzipped tarball named
\texttt{model.tar.gz}. The internal layout depends on the framework, but the
general structure for a PyTorch model is:

\begin{lstlisting}[style=bashcode]
model/
  code/
    inference.py        # entry point (see Step 3)
    requirements.txt    # extra dependencies (optional)
  model.pth             # the trained weights
\end{lstlisting}

Create the tarball:

\begin{lstlisting}[style=bashcode]
cd model/
tar -czvf model.tar.gz *
\end{lstlisting}

\begin{notebox}
For built-in algorithms (XGBoost, Linear Learner, etc.) SageMaker's own training
jobs produce a compatible \texttt{model.tar.gz} automatically --- you can skip
the custom packaging step.
\end{notebox}

% ---------------------------------
\subsection{Step 2 --- Upload the Artifact to S3}
% ---------------------------------

\begin{lstlisting}[style=bashcode]
aws s3 cp model.tar.gz s3://my-ml-artifacts/models/sentiment-v1/model.tar.gz
\end{lstlisting}

Note the S3 URI --- it will be referenced when creating the SageMaker Model
object. Use versioned paths (e.g., \texttt{sentiment-v1}, \texttt{sentiment-v2})
so rollbacks are trivial.

% ---------------------------------
\subsection{Step 3 --- Write the Inference Script}
% ---------------------------------

For custom models, SageMaker's framework containers (PyTorch, TensorFlow,
HuggingFace, scikit-learn) expose four overridable hooks in \texttt{inference.py}:

\begin{lstlisting}[style=pycode]
import json
import os
import torch

def model_fn(model_dir):
    """Load and return the model. Called once when the container starts."""
    model = torch.load(os.path.join(model_dir, "model.pth"),
                       map_location="cpu")
    model.eval()
    return model

def input_fn(request_body, request_content_type):
    """Parse the incoming request payload."""
    if request_content_type == "application/json":
        data = json.loads(request_body)
        return torch.tensor(data["inputs"], dtype=torch.float32)
    raise ValueError(f"Unsupported content type: {request_content_type}")

def predict_fn(input_data, model):
    """Run inference."""
    with torch.no_grad():
        return model(input_data)

def output_fn(prediction, response_content_type):
    """Serialize the prediction."""
    if response_content_type == "application/json":
        return json.dumps({"predictions": prediction.tolist()})
    raise ValueError(f"Unsupported content type: {response_content_type}")
\end{lstlisting}

\begin{tipbox}
If any of these four functions are missing, SageMaker uses sensible defaults
that handle JSON / NumPy payloads. For quick prototypes you can often ship with
just \texttt{model\_fn}.
\end{tipbox}

% ---------------------------------
\subsection{Step 4 --- Choose a Container Image}
% ---------------------------------

There are three options:

\begin{enumerate}[leftmargin=*]
    \item \textbf{Built-in SageMaker images} --- Maintained by AWS for PyTorch,
          TensorFlow, MXNet, scikit-learn, XGBoost, HuggingFace, etc. Resolve
          the URI via the SDK:
          \begin{lstlisting}[style=pycode]
from sagemaker import image_uris
image_uri = image_uris.retrieve(
    framework="pytorch",
    region="us-east-1",
    version="2.1.0",
    py_version="py310",
    image_scope="inference",
    instance_type="ml.m5.xlarge",
)
          \end{lstlisting}

    \item \textbf{Extend a SageMaker image} --- Start \texttt{FROM} an official
          image and add custom system packages or Python libraries. Push to ECR.

    \item \textbf{Bring Your Own Container (BYOC)} --- Fully custom Docker image
          that implements the SageMaker inference contract (exposes
          \texttt{/ping} and \texttt{/invocations} on port 8080). Push to ECR.
\end{enumerate}

\begin{warnbox}
If you use BYOC, the container must respond to \texttt{GET /ping} with
HTTP 200 within the timeout (default 60 s) or the endpoint will fail to come up.
This is the most common BYOC bug.
\end{warnbox}

% ---------------------------------
\subsection{Step 5 --- Create the SageMaker Model}
% ---------------------------------

Using the high-level SageMaker Python SDK:

\begin{lstlisting}[style=pycode]
from sagemaker.pytorch import PyTorchModel
import sagemaker

sess = sagemaker.Session()
role = "arn:aws:iam::123456789012:role/SageMakerExecutionRole"

model = PyTorchModel(
    model_data="s3://my-ml-artifacts/models/sentiment-v1/model.tar.gz",
    role=role,
    framework_version="2.1.0",
    py_version="py310",
    entry_point="inference.py",
    source_dir="model/code",          # where inference.py lives locally
    name="sentiment-model-v1",
    sagemaker_session=sess,
)
\end{lstlisting}

Behind the scenes this resolves the container image, uploads
\texttt{source\_dir} as a secondary artifact, and prepares a Model object ---
but does not yet create it in the account until \texttt{deploy()} is called.

% ---------------------------------
\subsection{Step 6 --- Deploy the Endpoint}
% ---------------------------------

A single call creates the EndpointConfig and Endpoint:

\begin{lstlisting}[style=pycode]
predictor = model.deploy(
    initial_instance_count=1,
    instance_type="ml.m5.xlarge",
    endpoint_name="sentiment-endpoint",
)
\end{lstlisting}

Provisioning typically takes \textbf{5--10 minutes}. The endpoint state
transitions: \texttt{Creating} $\rightarrow$ \texttt{InService}.

\paragraph{Common instance families:}
\begin{itemize}[leftmargin=*,nosep]
    \item \texttt{ml.t2.*, ml.m5.*} --- CPU, general purpose, cheap.
    \item \texttt{ml.c5.*} --- compute-optimized CPU, good for latency-sensitive
          classical ML.
    \item \texttt{ml.g4dn.*, ml.g5.*} --- NVIDIA T4 / A10G GPUs, cost-effective
          deep learning inference.
    \item \texttt{ml.p4d.*, ml.p5.*} --- A100 / H100 GPUs, LLMs and large
          vision models.
    \item \texttt{ml.inf1.*, ml.inf2.*} --- AWS Inferentia custom silicon,
          lowest cost-per-inference for supported models.
\end{itemize}

% ---------------------------------
\subsection{Step 7 --- Invoke the Endpoint}
% ---------------------------------

From the SDK:

\begin{lstlisting}[style=pycode]
result = predictor.predict({"inputs": [[0.1, 0.4, 0.9]]})
print(result)
\end{lstlisting}

Or from any language using \texttt{boto3} / AWS SDK:

\begin{lstlisting}[style=pycode]
import boto3, json

runtime = boto3.client("sagemaker-runtime", region_name="us-east-1")
resp = runtime.invoke_endpoint(
    EndpointName="sentiment-endpoint",
    ContentType="application/json",
    Body=json.dumps({"inputs": [[0.1, 0.4, 0.9]]}),
)
print(json.loads(resp["Body"].read()))
\end{lstlisting}

\begin{notebox}
SageMaker endpoints are \emph{private by default} --- they sit behind IAM, not
the public internet. To expose them externally, put an API Gateway + Lambda (or
direct API Gateway HTTP integration) in front, terminating auth there.
\end{notebox}

% ---------------------------------
\subsection{Step 8 --- Configure Autoscaling}
% ---------------------------------

Autoscaling is \emph{not} enabled by default. Attach a target-tracking policy
via Application Auto Scaling:

\begin{lstlisting}[style=pycode]
import boto3

asg = boto3.client("application-autoscaling")
resource_id = "endpoint/sentiment-endpoint/variant/AllTraffic"

asg.register_scalable_target(
    ServiceNamespace="sagemaker",
    ResourceId=resource_id,
    ScalableDimension="sagemaker:variant:DesiredInstanceCount",
    MinCapacity=1,
    MaxCapacity=8,
)

asg.put_scaling_policy(
    PolicyName="sentiment-scale",
    ServiceNamespace="sagemaker",
    ResourceId=resource_id,
    ScalableDimension="sagemaker:variant:DesiredInstanceCount",
    PolicyType="TargetTrackingScaling",
    TargetTrackingScalingPolicyConfiguration={
        "TargetValue": 70.0,
        "PredefinedMetricSpecification": {
            "PredefinedMetricType": "SageMakerVariantInvocationsPerInstance"
        },
        "ScaleInCooldown": 300,
        "ScaleOutCooldown": 60,
    },
)
\end{lstlisting}

% ---------------------------------
\subsection{Step 9 --- Monitor and Log}
% ---------------------------------

SageMaker emits metrics to CloudWatch automatically:

\begin{itemize}[leftmargin=*,nosep]
    \item \texttt{Invocations}, \texttt{InvocationsPerInstance}
    \item \texttt{ModelLatency}, \texttt{OverheadLatency}
    \item \texttt{Invocation4XXErrors}, \texttt{Invocation5XXErrors}
    \item \texttt{CPUUtilization}, \texttt{GPUUtilization}, \texttt{MemoryUtilization}
\end{itemize}

Container \texttt{stdout}/\texttt{stderr} flows to CloudWatch Logs under
\texttt{/aws/sagemaker/Endpoints/<endpoint-name>}.

For production health, also enable:
\begin{itemize}[leftmargin=*,nosep]
    \item \textbf{Model Monitor} --- detects data drift, feature skew, model
          quality drift, and bias over time.
    \item \textbf{Clarify} --- explainability and bias reports.
    \item \textbf{CloudWatch Alarms} --- page on \texttt{ModelLatency} p99 or
          5XX rate.
\end{itemize}

% ---------------------------------
\subsection{Step 10 --- Update, Roll Back, or Delete}
% ---------------------------------

\paragraph{Rolling update.} To ship a new model version, create a new Model
and a new EndpointConfig, then update the existing Endpoint:

\begin{lstlisting}[style=pycode]
sm = boto3.client("sagemaker")
sm.update_endpoint(
    EndpointName="sentiment-endpoint",
    EndpointConfigName="sentiment-config-v2",
)
\end{lstlisting}

SageMaker performs a \textbf{blue/green} swap: the new fleet is provisioned,
health-checked, then traffic is shifted. The old fleet is torn down only after
the new fleet is healthy.

\paragraph{Delete.} Clean up to stop incurring charges (real-time endpoints
bill per second regardless of traffic):

\begin{lstlisting}[style=pycode]
predictor.delete_endpoint()
predictor.delete_model()
\end{lstlisting}

\begin{warnbox}
A running \texttt{ml.g5.xlarge} endpoint costs roughly \$1+/hour even when
idle. \textbf{Always delete test endpoints} or use Serverless Inference for
development workloads.
\end{warnbox}

% =========================================================
\section{Deployment Variants}
% =========================================================

% ---------------------------------
\subsection{Serverless Inference}
% ---------------------------------

No instance management; pay per invocation and duration. Ideal for spiky or
low-traffic workloads.

\begin{lstlisting}[style=pycode]
from sagemaker.serverless import ServerlessInferenceConfig

serverless_config = ServerlessInferenceConfig(
    memory_size_in_mb=4096,
    max_concurrency=20,
)
predictor = model.deploy(serverless_inference_config=serverless_config)
\end{lstlisting}

\textbf{Trade-offs:} cold starts (hundreds of ms to seconds), 6 GB memory ceiling,
no GPU support.

% ---------------------------------
\subsection{Asynchronous Inference}
% ---------------------------------

Requests are queued in S3; responses returned via SNS notifications. Suitable
for large payloads (up to 1 GB) or long-running inferences (up to 1 hour).

\begin{lstlisting}[style=pycode]
from sagemaker.async_inference import AsyncInferenceConfig

async_config = AsyncInferenceConfig(
    output_path="s3://my-ml-artifacts/async-outputs/",
    max_concurrent_invocations_per_instance=4,
)
predictor = model.deploy(
    initial_instance_count=1,
    instance_type="ml.m5.xlarge",
    async_inference_config=async_config,
)
\end{lstlisting}

% ---------------------------------
\subsection{Batch Transform}
% ---------------------------------

One-shot bulk scoring of a dataset in S3, no persistent endpoint:

\begin{lstlisting}[style=pycode]
transformer = model.transformer(
    instance_count=2,
    instance_type="ml.m5.xlarge",
    output_path="s3://my-ml-artifacts/batch-outputs/",
)
transformer.transform(
    data="s3://my-ml-artifacts/batch-inputs/",
    content_type="application/json",
)
transformer.wait()
\end{lstlisting}

% ---------------------------------
\subsection{Multi-Model Endpoints (MME)}
% ---------------------------------

Host \emph{many} models behind one endpoint, loaded dynamically from S3 on
demand. Ideal when you have hundreds or thousands of small models (e.g., one
per customer) with sparse traffic.

% ---------------------------------
\subsection{Multi-Container Endpoints}
% ---------------------------------

Host multiple containers behind one endpoint. Invoke a specific container by
name, or chain them in a serial inference pipeline (e.g., preprocessing
$\rightarrow$ model $\rightarrow$ postprocessing).

% =========================================================
\section{Production Best Practices}
% =========================================================

\begin{enumerate}[leftmargin=*]
    \item \textbf{Infrastructure as Code} --- Define endpoints in CloudFormation,
          CDK, or Terraform. Avoid console-driven deployments.
    \item \textbf{CI/CD} --- Use SageMaker Pipelines or CodePipeline for
          train $\rightarrow$ evaluate $\rightarrow$ register $\rightarrow$
          deploy flows. Store approved models in the \emph{Model Registry}.
    \item \textbf{Blue/Green + Canary} --- Use \texttt{DeploymentConfig} with
          canary or linear traffic shifting for risk-averse rollouts with
          automatic CloudWatch-alarm-triggered rollback.
    \item \textbf{VPC isolation} --- Place endpoints in a VPC with no internet
          egress when handling sensitive data. Use VPC endpoints for S3 and ECR
          pulls.
    \item \textbf{Encryption} --- Enable KMS encryption for model artifacts in
          S3, endpoint EBS volumes, and inter-node traffic.
    \item \textbf{Right-sizing} --- Use SageMaker Inference Recommender to
          benchmark instance types against your latency and throughput SLOs.
    \item \textbf{Shadow testing} --- Deploy new variants with 0\% traffic and
          mirror production requests to validate before cutover.
    \item \textbf{Cost hygiene} --- Tag endpoints with project / environment /
          owner and review Cost Explorer weekly. Use Savings Plans for steady
          workloads.
\end{enumerate}

% =========================================================
\section{Troubleshooting Checklist}
% =========================================================

\begin{longtable}{p{5.5cm} p{9.2cm}}
\toprule
\textbf{Symptom} & \textbf{Likely Cause / Fix} \\
\midrule
Endpoint stuck in \texttt{Creating} then \texttt{Failed} &
Check CloudWatch logs. Usually \texttt{model\_fn} throws, container \texttt{/ping}
fails, or \texttt{requirements.txt} install times out. \\
\addlinespace
\texttt{AccessDenied} on S3 &
SageMaker execution role missing \texttt{s3:GetObject} on the artifact bucket. \\
\addlinespace
\texttt{ModelError: 413 Payload Too Large} &
Real-time endpoints cap payloads at 6 MB. Use Async Inference for larger inputs. \\
\addlinespace
\texttt{ModelLatency} spikes &
Enable autoscaling; increase instance size; batch requests; consider GPU or
Inferentia. \\
\addlinespace
Cold start on first invocation &
Pre-warm by invoking after deploy; use provisioned concurrency on serverless. \\
\addlinespace
High cost with low traffic &
Switch to Serverless Inference or Async Inference; delete idle endpoints. \\
\bottomrule
\end{longtable}

% =========================================================
\section{Summary}
% =========================================================

The canonical SageMaker real-time deployment flow is:

\begin{enumerate}[leftmargin=*,nosep]
    \item Package model $\rightarrow$ \texttt{model.tar.gz}.
    \item Upload to S3.
    \item Write \texttt{inference.py} (four hooks).
    \item Pick or build a container image.
    \item Create a \texttt{Model} (binds artifact + image + role).
    \item Create an \texttt{EndpointConfig} (instance type, count, scaling).
    \item Create / update an \texttt{Endpoint}.
    \item Invoke via \texttt{sagemaker-runtime}.
    \item Attach autoscaling and CloudWatch alarms.
    \item Monitor with Model Monitor; iterate with blue/green updates.
\end{enumerate}

The SDK collapses steps 5--7 into a single \texttt{model.deploy()} call, which
is why SageMaker feels lightweight in practice despite the rich underlying
primitives.

\end{document}
